% E5: How much of the dictionary is used, and how sparse is the code really?
% Drop-in section. Uses only what the main paper already loads.
% Occupancy values are filled from discovery run counterfactual-20260924-154113.
% Token-level L0 values are \RESULT placeholders pending a capture run.
% Artefact keys for each table are named in a comment above it.

\providecommand{\RESULT}[1]{\textbf{[#1]}}

\section{Dictionary Occupancy and the Sparsity of the Code}
\label{sec:dictionary}

The transcoders are described as sparse and are given a sixteen-fold expansion,
$d_{\mathrm{features}} = 16\,d_{\mathrm{model}} = 16384$ per layer, under an
$\ell_1$ penalty of $10^{-4}$. Neither the sparsity nor the occupancy of that
dictionary is reported anywhere. This section measures what a discovery run can
establish about both, and is as much about what such a run \emph{cannot}
establish, because the natural statistic is an upper bound that is loose by up
to a factor of fifty and the natural dead-feature count is not identifiable at
the sample sizes used so far.

\subsection{Hypothesis}

\paragraph{H5.} The dictionary is substantially used: the fraction of features
that never activate is small, and the number active at once is a small fraction
of the dictionary.

\emph{Falsified if} a large share of features is genuinely dead, or if the
active fraction is comparable to an unexpanded basis, either of which would mean
the expansion is not buying a sparse overcomplete code.

\emph{Decision quantities and the rule.} The never-fired fraction, reported as a
dead-feature count \emph{only} when the sample size can distinguish dead from
rare, in the sense of Section~\ref{sec:dict-identifiability}; and the mean
number of features active per token, which requires a measurement the discovery
pipeline does not make.

\subsection{Dead is not the same as unobserved}
\label{sec:dict-identifiability}

Counting features that never fired does not count dead features. It bounds them.
A feature that fires once in six hundred observations will usually be unseen in
a hundred, and the rate matters here because the paper's own case-study feature
$\mathtt{L11}\!:\!\tau\!=\!0.7\!:\!\mathtt{F9970}$ fires at $0.00166$, chosen for
study precisely because it is rare.

With no firings in $N$ observations, the exact one-sided upper confidence bound
on a feature's rate solves $(1-p)^N = \alpha$, so
\begin{equation}
p_{\max} = 1 - \alpha^{1/N},
\qquad
N_{\mathrm{needed}}(p) = \frac{\log \alpha}{\log(1-p)},
\label{eq:dict-bound}
\end{equation}
the familiar rule of three at $\alpha = 0.05$. A run separates dead from rare
only when $p_{\max}$ falls below the rate worth detecting. At the $160$
observations available here, $p_{\max} = 1.85\%$, which is eleven times the
reference rate; a dead-feature claim would need about $1804$ observations, and
the paper's own $13{,}835$-observation discovery run clears that with room to
spare. The occupancy below is therefore reported as a bound, and the dead count
is left to the larger run.

Two properties of the pipeline make even that bound optimistic, and both are
stated rather than corrected. Discovery observations are consecutive frames from
episodes, so a feature's firings are correlated across them and the effective
sample is smaller than the nominal one. And discovery reduces the code by a
maximum over the fifty action positions, so these are per-observation rates: a
feature counts as firing if it fired at any token.

\subsection{Results}

% Filled from dictionary_health.csv and dictionary_health.json -> layers, summary.
\begin{table}[t]
\caption{Dictionary occupancy by layer, from $160$ observations. Rates are
per-observation, over a maximum across the fifty action positions. The last
column is the expected number of features active in an observation, which
bounds per-token $L_0$ from above.}
\label{tab:dict-occupancy}
\begin{center}
\small
\setlength{\tabcolsep}{4.5pt}
\begin{tabular}{crrrrr|crrrrr}
\hline
$\ell$ & Never & \% & Med.\ rate & $\mathbb{E}[\mathrm{act}]$ & \% dict &
$\ell$ & Never & \% & Med.\ rate & $\mathbb{E}[\mathrm{act}]$ & \% dict \\
\hline
0 & $12956$ & $79$ & $0.99$ & $2715$ & $17$ & 9  & $84$   & $1$  & $0.57$ & $5258$ & $32$ \\
1 & $12154$ & $74$ & $0.94$ & $2875$ & $18$ & 10 & $226$  & $1$  & $0.44$ & $3817$ & $23$ \\
2 & $11461$ & $70$ & $0.96$ & $3483$ & $21$ & 11 & $75$   & $0$  & $0.44$ & $4977$ & $30$ \\
3 & $8833$  & $54$ & $0.96$ & $4360$ & $27$ & 12 & $249$  & $2$  & $0.42$ & $4733$ & $29$ \\
4 & $1093$  & $7$  & $0.82$ & $7446$ & $45$ & 13 & $623$  & $4$  & $0.33$ & $4106$ & $25$ \\
5 & $1150$  & $7$  & $0.69$ & $6132$ & $37$ & 14 & $1006$ & $6$  & $0.28$ & $3754$ & $23$ \\
6 & $218$   & $1$  & $0.70$ & $7233$ & $44$ & 15 & $877$  & $5$  & $0.26$ & $3977$ & $24$ \\
7 & $1395$  & $9$  & $0.43$ & $3576$ & $22$ & 16 & $4979$ & $30$ & $0.19$ & $2394$ & $15$ \\
8 & $131$   & $1$  & $0.58$ & $5762$ & $35$ & 17 & $9301$ & $57$ & $0.15$ & $1770$ & $11$ \\
\hline
\multicolumn{12}{l}{Never fired overall: $66{,}811$ of $294{,}912$ ($22.7\%$), per-layer range $0\%$ to $79\%$.} \\
\multicolumn{12}{l}{An unseen feature fires at most $1.85\%$ at this sample size, so none of these is a dead count.} \\
\hline
\end{tabular}
\end{center}
\end{table}

\paragraph{The dictionary is bimodal, and the shape is the finding.}
At the first four layers most of the dictionary is never touched while the
features that are touched are almost always on: at layer $0$, $79\%$ of features
never fire and the median rate among the rest is $0.99$. The middle of the stack
inverts this, with under $2\%$ unused at layers $6$ and $8$ through $12$ and
median rates falling toward $0.4$. The last two layers return to a sparsely
occupied dictionary. So the expansion is bought at very different prices along
the stack, and an aggregate expansion factor describes none of it well. Note
this profile is the near-mirror of the flow-time result of
Section~\ref{sec:tau}: the layers with the least-used dictionaries are also the
layers whose codes barely depend on $\tau$.

\paragraph{The sparsity number the paper needs is not in this measurement.}
The expected active count runs from $1770$ to $7446$ features, eleven to
forty-five per cent of the dictionary. Read literally that is not a sparse code.
But it is an upper bound, and a loose one: a feature counts as active if it fires
at any of the fifty action positions, so a code with a per-token $L_0$ of $100$
and no overlap between tokens would present here as $5000$ active. The bound is
loose by up to the token count. A per-token $L_0$ can be read directly from the
latents captured during a rollout, which the runtime already records, and until
it is reported the paper has no sparsity figure at all.

% Filled from the per-token L0 analysis over a rollout capture (transcoder_capture/events.jsonl).
\begin{table}[t]
\caption{Per-token sparsity from captured rollout latents. This is the statistic
the word ``sparse'' refers to; Table~\ref{tab:dict-occupancy} bounds it from
above by up to the number of action tokens.}
\label{tab:dict-l0}
\begin{center}
\small
\begin{tabular}{lcccc}
\hline
Quantity & Median & Mean & 90th pct & Max \\
\hline
$L_0$ per token & \RESULT{val} & \RESULT{val} & \RESULT{val} & \RESULT{val} \\
As \% of the $16384$ dictionary & \RESULT{val} & \RESULT{val} & \RESULT{val} & \RESULT{val} \\
\hline
\multicolumn{5}{l}{Range across layers: \RESULT{lo} to \RESULT{hi}; tokens measured: \RESULT{n}} \\
\hline
\end{tabular}
\end{center}
\end{table}

\subsection{Relation to the rest of the paper}

\paragraph{Architecture and training.} The held-out reconstruction loss of
$0.046$ establishes that the transcoders reproduce the MLPs; it says nothing
about the sparsity that is the other half of their objective. These two tables
are the missing half. The occupancy result also gives the expansion factor an
empirical footing it currently lacks: a sixteen-fold expansion is generous at
layer $0$, where four fifths of it goes unused, and tight in the middle of the
stack, where nearly all of it is live.

\paragraph{Feature discovery.} The interesting score divides a top-$K$ mean by a
standard deviation, which rewards features that are rare and strong. Occupancy
explains where such features can be found: at layers $0$ to $3$ the live
features fire almost always and so score poorly by construction, while the
sparsely firing features that the ranking favours live in the later layers. The
paper's own top-ranked examples sit at layers $11$, $12$ and $17$, which is
consistent with this and worth stating rather than leaving as coincidence.

\paragraph{Circuit construction.} A layer where thousands of features are active
in every observation is a layer where an attribution frontier has thousands of
candidate parents. Taken with the flow-time redundancy of Section~\ref{sec:tau},
this is the second structural reason a backward trace inflates.

\subsection{Limitations and what would settle it}

The occupancy figures come from one $160$-observation run and are bounds, not
counts; Eq.~\ref{eq:dict-bound} gives the sample size that would convert them.
The active-count column bounds per-token $L_0$ from above and should not be
quoted as sparsity. Three measurements would close the section: the same
occupancy analysis on the $13{,}835$-observation discovery set, where dead
becomes identifiable; per-token $L_0$ from a rollout capture; and a small sweep
of the $\ell_1$ coefficient to place the operating point on a
reconstruction-against-sparsity frontier rather than asserting it. Only the last
needs new training.
